\subsection{The Cycle Finding Problem}

\frame{

\frametitle{\secname: \subsecname}

\begin{columns}
\column{0.5\textwidth}

\begin{itemize}
\item<1-> Take some function $f:[n]\to[n]$,\\and some starting value $x$.
\begin{itemize}
\item<1-> Notation: $[n]$ is the numbers $\{1,\dots,n\}$, and $f$ is a function that takes a number from $[n]$, and outputs a number in $[n]$
\end{itemize}
\item<2-> By pidgeonhole principle, the sequence\\
$x, f(x), f(f(x)),\dots,f^i(x),\dots$\\
will repeat
\item<3-> Find $s,t$ such that $f^s(x)=f^{s+t}(x)$
\end{itemize}

\column{0.5\textwidth}
\visible<2->{
\begin{figure}
\centering
\begin{tikzpicture}[every node/.style={node distance=1cm}]
\node (x0) {$x$};
\node (x1)[left=of x0] {$f(x)$};
\node (x2)[below left=of x1] {$f(f(x))$};
\node (s0)[below=of x2] {$f^s(x)$};
\node (s1)[below right=of s0] {$f^{s+1}(x)$};
\node (tm2)[above right=of s1] {$f^{t-2}(x)$};
\node (tm1)[above left=of tm2] {$f^{t-1}(x)$};

\draw[-{Latex[length=8pt]}] (x0) -- (x1);
\draw[-{Latex[length=8pt]}] (x1) -- (x2);
\draw[-{Latex[length=8pt]}] (x2) edge[dotted] (s0);
\draw[-{Latex[length=8pt]}] (s0) -- (s1);
\draw[-{Latex[length=8pt]}] (s1) edge[dotted] (tm2);
\draw[-{Latex[length=8pt]}] (tm2) -- (tm1);
\draw[-{Latex[length=8pt]}] (tm1) -- (s0);
\end{tikzpicture}
\end{figure}
}

\end{columns}

}
